\scanpage{9}\mspage{8.}

donc

\[
\begin{aligned}
D^2&=\sum_{i,j}\underbrace{(e_i\otimes f_i+f_i\otimes e_i)(e_{\unsure{j|g}}\otimes f_j+f_j\otimes e_j)}_{-[e_i f_j\otimes f_i e_j+f_i e_j\otimes e_i f_j]}\\
&=-\sum_i(\eta\otimes f_i e_i-\eta\otimes e_i f_i)=2\eta\sum e_i f_i
\end{aligned}
\]

\[
\text{donc}\qquad\tfrac12D^2=\eta\otimes\sum e_i f_i
\]

et par suite

\[
\varphi(\delta)=\sum e_i\wedge f_i\in H^2(B)
\]

\uline{C'est la classe de polarisation canonique} de la \uline{jacobienne $B$}. D'ailleurs $\chi(\delta)=1-g+d$

% slanted
\margin{vérifier que le signe est correct \unsure{ici}}

\deleted{\unsure{R|et}} Utilisons la polarisation canonique \inserted{par le diviseur $\Theta$} de $B$, définissant un isom can

\[
c:B\longrightarrow A
\]

et transformons les $A$-torseurs de la fin de d) par $c^{-1}$. \uline{On sait que} $\underline{\mathrm{Alb}}^1_{X/S}$ \uline{est transformé} en $\underline{\mathrm{Pic}}^{-1}_{X/S}=P^{-1}$ (\uline{attention au signe} !) donc $\underline{\mathrm{Alb}}^{\chi(\delta)}_{X/S}$ devient $P_{\deleted{\unsure{X/S}}}^{-\chi(\delta)}=P^{g-1-d}$, d'autre part $\underline{\mathrm{Pic}}^{\varphi(\delta)=\text{pol. can. de } A}_{B/S}$ \deleted{canoniquement} \inserted{devient} $\deleted{\unsure{\mathrm{Pic}^{g-1}|\underline{\mathrm{Pic}}^1}}\ P^{g-1}$ (isom. canonique respectant les div. $\Theta$), enfin $P^\delta\overset B\times(A,\underset{\substack{\shortparallel\\c}}{\widetilde{\varphi(\delta)}})$ devient $P^\delta=P^d$. Donc le torseur envisagé, produit \inserted{contracté} des trois précédents, est $P^{2g-2}$. On trouve une section canonique de ce dernier : ce n'est pas très excitant, car on s'attend à trouver la section canonique $K$ (à prouver !).
